\chapter*{Conclusiones y recomendaciones}
\addcontentsline{toc}{chapter}{Conclusiones y recomendaciones}
\markboth{Conclusiones y recomendaciones}{Conclusiones y recomendaciones}

\section*{Conclusiones}

El presente trabajo cumplió con el objetivo general planteado: se desarrolló \emph{EduFEM}, un software educativo de escritorio escrito en el lenguaje de programación Python sobre bibliotecas numéricas de código abierto, para el análisis por el Método de los Elementos Finitos (MEF) de problemas bidimensionales de elasticidad lineal en tensión plana y deformación plana, con elementos cuadriláteros isoparamétricos Q4 y Q9. La herramienta integra, en un único entorno coherente, las tres dimensiones que se propusieron como ejes del proyecto: la formulación matemática del método, la construcción y visualización interactiva del modelo, y la verificación numérica de los resultados. El trabajo regresa así a la doble brecha que lo originó —una dificultad de aprendizaje genuina, ligada a la naturaleza abstracta del método, y la insuficiencia de herramientas que la atiendan exponiendo el proceso de cálculo y la respuesta de manera transparente y guiada— y a la interrogante que planteó como problema: de qué manera el uso del software como caja negra afecta el criterio del estudiante para interpretar la respuesta estructural. EduFEM la atiende ofreciendo un entorno que conserva el rigor numérico —verificado y validado— a la vez que expone, sobre el modelo concreto del estudiante, el procedimiento que produce la respuesta y los medios para juzgarla.

A continuación se presentan las conclusiones específicas, ordenadas en correspondencia con los objetivos planteados y con el capítulo que desarrolla cada uno.

En relación con el diagnóstico de la brecha (\autoref{cap:marco_teorico}, \autoref{sec:antecedentes}), la revisión de los antecedentes y el análisis comparativo de la \autoref{tab:comparativa} establecieron que los recursos con que el estudiante aprende el MEF le devuelven la respuesta estructural sin el procedimiento que la produce ni los medios para juzgarla: el software comercial por diseño, y el antecedente educativo más directo, ED-Elas2D, por reconocimiento de sus propios autores, que describen su post-proceso como una herramienta meramente descriptiva que no explica el suavizado nodal. De ese diagnóstico salieron los atributos que la herramienta debía tener: exposición interactiva de cada matriz sobre el elemento, memoria de cálculo trazable al modelo, respuesta expuesta con su procedimiento y sus verificaciones, y una batería de verificación y validación publicada con el software.

En relación con la fundamentación teórica (\autoref{cap:marco_teorico}), se estableció la base que el motor, los módulos educativos y la memoria de cálculo comparten: la formulación del MEF para la elasticidad lineal plana con elementos isoparamétricos Q4 y Q9 —de la formulación variacional a la recuperación de tensiones—, la calidad de malla, los fenómenos numéricos y los criterios de verificación y validación, junto con los cuatro principios de aprendizaje que guían el diseño, apoyados en la enseñanza del método y en la investigación general sobre simulaciones, compromiso cognitivo y ejemplos resueltos. Esa base no es un resultado medido sino el fundamento del trabajo; lo que los resultados comprueban es que el motor construido sobre ella se comporta como la teoría predice, que las mismas expresiones que el capítulo expone son las que el código implementa, los módulos exhiben y la memoria sustituye numéricamente, y que cada decisión de diseño remite a un principio documentado.

En relación con el desarrollo del software (\autoref{cap:diseno}; resultado en la \autoref{sec:cobertura-canal}), se construyó un núcleo numérico completo en \texttt{NumPy} y \texttt{SciPy} que cubre la totalidad del canal de cálculo del MEF para elementos isoparamétricos —funciones de forma bilineales y bicuadráticas, Jacobiano, matriz $\bm{B}$, matriz $\bm{D}$ para ambos estados planos, rigidez elemental por cuadratura de Gauss, ensamblaje global, imposición de restricciones incluidos los desplazamientos prescritos no homogéneos, y recuperación de tensiones por extrapolación desde los puntos de Gauss—, con un esquema de indexación de grados de libertad robusto frente a identificadores de nodo no contiguos. Sobre él se logró una aplicación organizada en las tres fases canónicas del análisis —pre-proceso, proceso y post-proceso— que comparten un único lienzo interactivo, con selección bidireccional entre las tablas y el lienzo, deshacer y rehacer sobre el modelo completo y validación de la salud del modelo previa a la resolución; ocho módulos educativos que exponen, paso a paso y sobre el modelo real que el estudiante ha construido, las etapas del canal desde el mapeo isoparamétrico hasta el ensamblaje global (M1 a M7, \autoref{tab:modulos}), más la calidad de malla en el pre-proceso (M0); un post-proceso que expone la solución y la recuperación de tensiones con sus modos crudo y suavizado, la sonda con círculo de Mohr, la vista tridimensional, las reacciones y la tabla de resultados; y una memoria de cálculo en PDF que reconstruye el procedimiento paso a paso con las mismas funciones del motor, generada sin error para Q4 y Q9 y reproducida en el \autoref{anexo:memoria}, con intercambio de datos por DXF y CSV verificado por ida y vuelta e idempotencia. A diferencia de las visualizaciones genéricas de los textos, los módulos operan directamente sobre el elemento seleccionado en el lienzo, conectando la abstracción matemática con la geometría concreta del problema: es la decisión de diseño que materializa el aporte pedagógico central del trabajo.

En relación con la verificación y validación del motor (\autoref{cap:resultados}), se aplicó una batería de tres niveles con criterios de aceptación fijados a priori en los propios guiones. La verificación de código por el método de soluciones manufacturadas, en tensión plana y en deformación plana y sobre mallas regulares y distorsionadas, confirmó las tasas de convergencia asintóticas teóricas $\mathcal{O}(h^{p+1})$ en norma $L^2$ y $\mathcal{O}(h^{p})$ en seminorma $H^1$ —$\mathcal{O}(h^{2{,}00})$ y $\mathcal{O}(h^{1{,}00})$ para Q4, $\mathcal{O}(h^{3{,}00})$ y $\mathcal{O}(h^{2{,}00})$ para Q9—, y el campo de tensiones recuperado converge a $\mathcal{O}(h^{1{,}5})$ en Q4 y $\mathcal{O}(h^{2})$ en Q9, con una matriz de extrapolación que reproduce exactamente los campos polinómicos de cada elemento; esas tasas son evidencia rigurosa de que el motor está libre de errores de orden. La contrastación con referencias externas se realizó con la viga de Timoshenko: frente a la solución analítica de Timoshenko y Goodier los errores fueron del $0{,}04\,\%$ en la tensión normal $\sigma_x$, del $0{,}26\,\%$ en la flecha y de hasta el $2{,}9\,\%$ en la tensión cortante; frente a un modelo equivalente en SAP2000, del $0{,}21\,\%$ en $\sigma_x$ y de hasta el $0{,}56\,\%$ en desplazamientos; y la suma de las reacciones equilibra la carga total a precisión de máquina. La membrana de Cook reprodujo el caso de referencia clásico y permitió ilustrar empíricamente el bloqueo por cortante (\emph{shear-locking}) del elemento Q4 frente a la convergencia rápida del Q9: para igual número de grados de libertad (GDL), el Q9 alcanzó un error de $-0{,}044\,\%$ donde el Q4 conservaba aún un $-0{,}594\,\%$. La contrastación se apoya en soluciones analíticas y en un modelo equivalente en un programa comercial, según la acepción declarada en el \autoref{cap:marco_teorico}; la validación en sentido estricto, contra mediciones físicas, queda fuera del alcance de este trabajo.

Por último, en relación con la validación del contenido y la coherencia del software (\autoref{cap:resultados}, \autoref{sec:contenido-didactico}), la tabla de especificaciones estableció que las quince destrezas de interpretación y verificación de la respuesta y los tres conceptos del canal que un consenso publicado de 67 expertos de la industria y de la academia exige al egresado tienen instrumento identificable en EduFEM, y que las de ranking más alto en ese consenso —mallado y verificación y validación— son las que el software ejercita sobre casos de referencia y con la memoria de cálculo. La matriz de principios estableció que los cuatro principios de aprendizaje tienen una encarnación comprobable en el diseño. Los dos instrumentos, construidos por vías distintas, señalan el mismo conjunto de decisiones —conversión de elemento y refinamiento, modo crudo y suavizado, equilibrio y diagnóstico en la memoria— como el corazón didáctico de la herramienta, y dejan a la vista sus dos límites: el software individual de escritorio no sostiene el modo interactivo entre pares, y las destrezas de mallado se ejercitan sobre secuencias que el estudiante debe construir sin un generador automático.

En conjunto, los resultados confirman que EduFEM satisface los cinco objetivos específicos en los términos en que fueron enunciados —los dos primeros como diagnóstico y fundamento, no como resultados medidos; el tercero con módulos hasta el ensamblaje y las dos últimas etapas expuestas en el post-proceso y la memoria— y, con ello, el objetivo general. Confirman también, cláusula por cláusula, la parte de la hipótesis que esta etapa contrasta: el software expone cada etapa del canal y la respuesta con su procedimiento, su motor reproduce las tasas teóricas con errores acotados frente a las referencias y evidencia el bloqueo por cortante del Q4 frente al Q9, y su contenido cubre las destrezas que los expertos exigen y encarna los principios que la literatura documenta. Que esa exposición mejore el criterio del estudiante para interpretar la respuesta estructural es la parte de la hipótesis que se fundamenta por analogía con las herramientas de la misma familia que sí midieron u observaron ese efecto, y que la prueba de campo de las recomendaciones contrastará.

\section*{Aportes del trabajo}

El principal aporte de este trabajo es una herramienta de software libre, funcional y verificada, que articula tres capas habitualmente disociadas en los recursos educativos sobre el MEF: un motor de cálculo riguroso, una interfaz de modelado interactiva y un conjunto de módulos que explican el método sobre el modelo real del estudiante.

Desde el punto de vista numérico, el motor implementa de manera completa y verificada la formulación isoparamétrica Q4 y Q9 para los dos estados planos de la elasticidad, con un esquema de indexación de GDL robusto frente a identificadores de nodo no contiguos —que desacopla la identidad del nodo de su índice de ensamblaje— y con tratamiento de cargas volumétricas, desplazamientos prescritos no homogéneos y normas de error $L^2$ y $H^1$ integradas con cuadratura de orden adecuado. La validación acota el error en las magnitudes primarias por debajo del $0{,}3\,\%$ frente a la solución analítica y del $0{,}6\,\%$ frente al modelo equivalente en SAP2000, para los problemas de referencia estudiados.

Desde el punto de vista pedagógico, el aporte distintivo es la decisión de diseño de que cada módulo educativo opere sobre el elemento que el estudiante selecciona en el lienzo compartido, en lugar de sobre un ejemplo prefabricado. Esto convierte la exploración del Jacobiano, de la matriz $\bm{B}$ o de la cuadratura de Gauss en una experiencia ligada al problema concreto, y permite que fenómenos sutiles como el bloqueo por cortante o los modos espurios (\emph{hourglass}) se observen en el modelo propio. La memoria de cálculo en PDF eleva esta capa a un segundo aporte pedagógico: genera, sobre el modelo concreto del usuario, un documento que reconstruye el procedimiento completo con sustitución numérica paso a paso —de las funciones de forma a las tensiones nodales— y lo cierra con un diagnóstico de equilibrio y condicionamiento. Frente a las visualizaciones interactivas, que muestran \emph{cómo} se calcula, la memoria deja un registro escrito que el alumno puede seguir, reproducir por sus propios medios y usar para verificar la respuesta, cerrando el lazo entre la teoría, la exploración y el cálculo manual. Frente al antecedente más directo, cuyo post-proceso devuelve la respuesta sin explicarla, EduFEM la expone con sus modos crudo y suavizado, sus reacciones contrastadas con la carga y su círculo de Mohr: es la parte del diseño que atiende de forma directa al criterio sobre la respuesta.

Desde el punto de vista metodológico, el trabajo aporta dos instrumentos reutilizables para evaluar recursos de enseñanza del MEF sin panel propio de expertos: la tabla de especificaciones construida sobre los ítems del consenso publicado de Pérez-Santiago y Campos, y la matriz de principios. Sus celdas fijan, además, qué debe observar la prueba de campo.

\section*{Limitaciones}

Las fronteras del trabajo se fijaron a priori en la Introducción (\emph{Alcance y limitaciones}): elasticidad lineal estática en dos dimensiones —sin sistemas estructurales discretos, placas, cáscaras ni problemas tridimensionales—, elementos cuadriláteros Q4 y Q9, material elástico lineal e isótropo, malla construida manualmente o por generación estructurada, y validación de la dimensión educativa limitada al contenido y a la coherencia del diseño. A ellas se suman las que el desarrollo y los resultados pusieron de manifiesto. En primer lugar, validez de contenido no es efectividad: los instrumentos establecen que el software expone lo que el criterio sobre la respuesta exige, no que el estudiante lo adquiera, y las tablas las construye el autor sobre un criterio externo publicado, no un panel independiente. En segundo lugar, la cobertura del canal por módulos interactivos llega hasta el ensamblaje: la solución del sistema y la recuperación de tensiones se exponen en el post-proceso y en la memoria, y no en un módulo propio; y las destrezas de mallado se ejercitan sobre secuencias de mallas que el estudiante debe construir a mano. En tercer lugar, el elemento Q4 exhibe bloqueo por cortante bajo flexión dominante, fenómeno que el software ilustra deliberadamente con fines didácticos pero que no mitiga mediante técnicas correctivas, y en deformación plana la matriz constitutiva degenera hacia la incompresibilidad cuando $\nu \to 0{,}5$, régimen que el módulo M4 señala pero que el software no resuelve. En cuarto lugar, el solucionador emplea una factorización directa sobre la matriz de rigidez almacenada en formato disperso, esquema eficiente para los tamaños de modelo propios de un contexto educativo ---e incluso para mallas de decenas de miles de grados de libertad, según cuantifica el \autoref{cap:resultados}--- pero cuyo coste crece de forma desfavorable en el extremo de mallas masivas. Finalmente, la validación descansa en tres casos de referencia y en un único caso con material de ingeniería civil; una batería más amplia reforzaría la confianza ante geometrías y cargas más diversas, y la carga superficial linealmente variable no interviene en ningún caso de referencia.

\section*{Recomendaciones y trabajo futuro}

A partir de las limitaciones anteriores se plantean tres líneas de continuación, una por capítulo, cada una con sus ítems concretos.

\subsection*{Sobre la formulación (\autoref{cap:marco_teorico})}

\textbf{Tratamiento del bloqueo.} Para abordar el bloqueo (\autoref{sec:locking}) —por cortante en flexión y volumétrico cuando $\nu \to 0{,}5$ en deformación plana— se recomienda incorporar técnicas correctivas como la integración reducida selectiva (SRI) o la formulación $\bar{\bm{B}}$ (B-bar) \autocite{hughes2000fem,bathe2014fem}; conviene recordar que el Q9 con integración completa tampoco está libre del bloqueo volumétrico, de modo que estas técnicas beneficiarían a ambos elementos. Más allá de su valor numérico, estas técnicas tienen un alto valor pedagógico: permitirían un módulo educativo dedicado que contraste, sobre un mismo modelo, el comportamiento bloqueado y el corregido, profundizando en la comprensión del fenómeno que la herramienta ya ilustra con el caso de referencia de Cook.

\textbf{Más tipos de elemento.} La biblioteca de elementos puede ampliarse con elementos triangulares (lineales y cuadráticos) y con cuadriláteros de transición, ofreciendo así una variedad que enriquezca el estudio comparativo de la convergencia y del comportamiento bajo distorsión de malla \autocite{onate2009structural,reddy2006introduction}.

\subsection*{Sobre el software (\autoref{cap:diseno})}

\textbf{Mallado automático.} Un generador de mallas sobre un contorno dibujado en el lienzo, con las métricas de calidad del módulo M0 como control, reduciría el trabajo manual que hoy exige un modelo de más de unas decenas de elementos y convertiría las destrezas de mallado de la tabla de especificaciones —selección del tamaño, refinamiento en zonas de concentración, prueba de convergencia— en operaciones de la propia herramienta.

\textbf{Solucionador para problemas de gran porte.} El motor ya ensambla la matriz de rigidez global vectorizada por lotes, la almacena en formato disperso comprimido (CSR) y resuelve el sistema reducido con una factorización LU dispersa directa con ordenamiento de mínimo grado, cuyo efecto cuantifica el \autoref{cap:resultados}. Las extensiones de mayor impacto sobre la escalabilidad consisten en incorporar una factorización de Cholesky —aplicable porque $\bm{K}_{ff}$ es simétrica y definida positiva (\autoref{sec:resolucion-tensiones})— y en ofrecer un solucionador iterativo precondicionado para problemas masivos. Estas mejoras preservarían la interfaz pública del solucionador y la versión legible del ensamblaje reservada a los módulos educativos, de modo que no obligarían a reescribir las capas de post-proceso ni dichos módulos.

\textbf{Extensión a tres dimensiones.} A más largo plazo, la extensión del motor a problemas tridimensionales con elementos hexaédricos (H8, H20) y tetraédricos constituiría un salto natural, aprovechando que gran parte de la formulación isoparamétrica y de la infraestructura de cuadratura es generalizable a tres dimensiones.

\subsection*{Sobre la evidencia (\autoref{cap:resultados})}

\textbf{Ampliación de la batería de validación.} La evidencia del \autoref{cap:resultados} descansa en tres casos de referencia y en un único caso con material de ingeniería civil, la viga de hormigón. Se recomienda incorporar un caso con carga superficial linealmente variable, que ningún caso de referencia ejercita, y al menos un caso civil en deformación plana con solución de referencia publicada —una presa de gravedad, un muro de contención o un túnel somero— con peso propio, de modo que la validación externa cubra el estado plano que la práctica civil emplea con más frecuencia.

\textbf{Prueba de campo del efecto sobre el criterio.} La línea de continuación más importante para consolidar la finalidad del proyecto es la prueba de campo que contraste la parte de la hipótesis que esta etapa fundamenta: que la exposición del canal y de la respuesta mejore el criterio del estudiante para interpretarla. Este trabajo deja ese diseño especificado en sus elementos esenciales. La población es la de la \autoref{sec:poblacion}, los estudiantes de la Carrera de Ingeniería Civil de la Universidad Autónoma Tomás Frías que cursan la asignatura en que se introduce el MEF, y la muestra, dirigida, un curso completo de esa asignatura. El diseño es de preprueba y posprueba con un solo grupo, o con dos cohortes si el calendario académico lo permite, en la línea de la comparación entre cohortes con que Bishay midió el efecto de sus herramientas \autocite[pp.~4-5, 13]{bishay2020teaching}. El tratamiento reproduce el uso previsto de la herramienta: construir el ejemplo canónico, validar su salud, resolverlo, recorrer los módulos M1 a M7 sobre un elemento propio, generar la memoria y reproducir a mano una etapa, y convertir el modelo de Q4 a Q9 y refinar la malla sobre la membrana de Cook. El instrumento mide el criterio de forma directa: un conjunto de resultados de elementos finitos para juzgar —una malla admisible y otra distorsionada, una deformada plausible y otra con una restricción faltante, tensiones nodales que cierran y otras que no, reacciones que equilibran la carga y otras que no, un Q4 bloqueado frente a un Q9 convergido—, con ítems construidos sobre las destrezas de la tabla de especificaciones, complementado por un cuestionario de percepción en la línea de la evaluación cualitativa que Lee aplicó en sus cursos \autocite[p.~168]{lee2015interactive}. El análisis compara la preprueba y la posprueba con estadística no paramétrica para muestras pareadas, adecuada al tamaño de un curso, y los archivos de proyecto y las memorias que cada alumno produce sirven de evidencia complementaria de su trabajo. Los resultados de esa prueba permitirían afinar el diseño de los módulos educativos sobre evidencia y aportarían el sustento empírico a la utilidad didáctica que el presente trabajo fundamenta por diseño y por la literatura.

En conjunto, estas líneas trazan un camino de evolución que preserva la identidad educativa de EduFEM al tiempo que amplía su alcance numérico, su biblioteca de elementos y, sobre todo, su fundamentación pedagógica empírica.
